\chapter{Composite actions}

Composite actions bundle multiple steps into a reusable action.
They live in your repository under \code{.github/actions/}.

\section{When to use composite actions}
Use them when you repeat the same steps across multiple workflows.
If the logic is very specific to one repo, a composite action is lighter than a reusable workflow.

\section{Action structure}
A composite action requires an \code{action.yml} with inputs and steps.

\begin{lstlisting}[language=YAML]
name: "setup tooling"
description: "Install tools used across workflows"
inputs:
  node-version:
    required: true
    default: "20"
runs:
  using: "composite"
  steps:
    - uses: actions/setup-node@v4
      with:
        node-version: ${{ inputs.node-version }}
    - run: npm ci
      shell: bash
\end{lstlisting}

\section{Using a composite action}
\begin{lstlisting}[language=YAML]
- uses: ./.github/actions/setup-tooling
  with:
    node-version: "20"
\end{lstlisting}

\section{Inputs and outputs}
Composite actions can expose outputs just like workflow steps.
Use them to pass versions or paths back to the caller.

\begin{lstlisting}[language=YAML]
outputs:
  cache-hit:
    value: ${{ steps.cache.outputs.cache-hit }}
\end{lstlisting}

\section{Limitations}
\begin{itemize}[nosep]
  \item Composite actions cannot define jobs or services.
  \item They run inside the caller job, so they share its permissions.
  \item Secrets must be passed as inputs explicitly.
\end{itemize}

\section{Versioning}
Tag composite actions when they are used across repositories.
If they are used only in one repo, keep them on the default branch and version by commit.

\section{Documentation}
Add a short README next to the action with inputs and outputs.
This makes it easier to reuse without reading the YAML.

\begin{patternbox}{Pattern: keep composite actions small}
A composite action should fit on one screen.
If it grows into multiple jobs or environments, use a reusable workflow.
\end{patternbox}

\begin{checklistbox}{Checklist: composite actions}
\begin{itemize}[nosep]
  \item Make inputs explicit and validated.
  \item Keep the action deterministic.
  \item Prefer a composite action over copy-paste.
\end{itemize}
\end{checklistbox}
